\index{covariance stationary}
\medskip
%\specsec{Definition:}
\definition{def:covstat}  A stochastic process $\{x_t\}$ is said to
be {\it covariance stationary\/} if it
satisfies the following two properties: (a) the mean is
independent of time,
 $ E x_t = E x_0$ for all $t$,
 and (b) the sequence of autocovariance matrices
 $E(x_{t+j} - E x_{t+j})(x_t - E x_t)'$
depends on the separation between dates
 $j = 0, \pm 1, \pm 2, \ldots$, but not on $t$.
\enddefinition
\smallskip





=================================

 We shall deduce
starting values for the mean and covariance
that make the process covariance
stationary, though our formulas are also useful for describing
what happens when we start from other initial conditions that
generate transient behavior that stops the process from being covariance
stationary.

=====================


We use
\definition{stable}
A square real valued matrix $A$ is said to be {\it stable\/} if
all of its eigenvalues modulus   are strictly less than
unity.
\enddefinition
\smallskip

We will make an assumption that guarantees that there exists an initial
condition $(\mu_0, \Sigma_0) = (Ex_0, E(x - Ex_0) (x-Ex_0)')$  that makes
the $x_t$ process covariance stationary.
Either of the following conditions works:

\medskip
\specsec{Condition A1:}  All of the eigenvalues of $A$ in
\Ep{statesp1} are strictly less than $1$ in modulus.




==============================


============================

A fixed point of the matrix difference equation \Ep{Sigma_diff} evidently  satisfies
$$ \Sigma_\infty = A \Sigma_\infty A' + C C' .  \EQN diff5 $$
A fixed point $ \Sigma_\infty  $
 is the covariance matrix $E (x_t -\mu) (x_t - \mu)'$ under  a stationary distribution of $x$.
%
% Thus, to compute $C_x(0)$, we  must solve
% $$ C_x(0)   =
%    A C_x(0) A^\prime  + C C' ,\EQN diff5 $$
% where $C_x(0) \equiv E (x_t -\mu) (x_t - \mu)'$.
Equation \Ep{diff5} is
a {\it discrete Lyapunov} equation in the $n \times n$ matrix
$\Sigma_\infty$.  It can be solved with the Matlab program {\tt
doublej.m}. % ^^|doublej.m|
\mtlb{doublej.m}%

============================


In the special case that $\Sigma_t = \Sigma_\infty$ that solves the discrete Lyapunov equation
\Ep{diff5}, $\Sigma_{t+j,t} = A^j \Sigma_\infty$ and so depends only on the gap $j$ between time
periods. In this case, an autocovariance  matrix sequence $\{\Sigma_{t+j,t}\}_{j=0}^\infty$ is often  also called an {\it
autocovariogram}.  \index{autocovariogram}%


===============================



diff2   mu motion

Sigma_diff sigma motion

============== Jan 29 2019 removed 


\section{Population regression}\label{sec:populationregression}%
\index{population regression}%
This section explains the notion of a population regression
equation.  \index{regression!population}%
  Suppose that we have a state-space system \Ep{statesp1} with
initial conditions that make it covariance stationary. We can use
the preceding formulas to compute the second moments of any pair
of random variables.  These moments let us compute a linear
regression.   Thus, let $X$ be a $p \times 1$ vector of random
variables somehow selected from the stochastic process $\{y_t\}$
governed by the system \Ep{statesp1}.
   For example,  let $p=  2m$, where
$y_t$ is an $m \times 1$ vector, and take $X =  \bmatrix{y_t \cr y_{t-1}\cr}$ for any $t \geq 1$.   Let $Y$ be any scalar
random variable selected from the $m \times 1$ stochastic process
$\{y_t\}$.  For example, take $Y = y_{t+1,1}$ for the same $t$
used to define $X$, where $y_{t+1,1}$ is the first component of
$y_{t+1}$.

  We consider the following least-squares approximation problem:
find a  $1\times p$ vector of real numbers $\beta$ that attain
$$ \min_\beta E (Y -  \beta X)^2.  \EQN leastsq  $$
Here $ \beta X$ is being  used to estimate $Y$, and we want the
value of $\beta$ that minimizes the expected squared error. The
first-order necessary condition for minimizing $E (Y -  \beta
X)^2$ with respect to $\beta$ is
$$ E  (Y -  \beta  X) X' = 0 , \EQN normal$$
which  can be rearranged as\NFootnote{That
 $E X'X$ is nonnegative definite implies that
the second-order conditions for a minimum of condition \Ep{leastsq} are
satisfied.}
$$ \beta =  (  E    YX') [E(X X')]^{-1}.  \EQN leastsq2  $$

  By using the formulas \Ep{diff2}, \Ep{diff5}, \Ep{diff6},  and
\Ep{ydiff2}, we can compute $E  X X'$ and $E Y X'$ for whatever
selection of $X$ and $Y$ we choose. The condition \Ep{normal} is
called the least-squares normal equation.  It states that the
projection error $Y - \beta X $ is orthogonal to $X$. Therefore,
we can represent $Y$ as
$$ Y =  \beta X + \epsilon \EQN regress $$
where $E  \epsilon X' =0$.   Equation \Ep{regress} is called a
population regression equation, and $\beta X $ is called the least-squares
projection of $Y$ on $X$ or the least-squares regression of $Y$ on
$X$. The vector $\beta$ is called the population least-squares
regression vector.
 The law of large numbers for continuous-state Markov processes,
\Theorem{LLNMarkovcontinuous},
\index{law of large numbers}%
states conditions  that guarantee that sample moments converge to
population moments, that is, ${1 \over S} \sum_{s=1}^S  X_s X_s'
\rightarrow E  X X' $ and ${ 1\over S} \sum_{s=1}^S  Y_s X_s'
\rightarrow E  Y X'$. Under those conditions, sample least-squares
estimates converge to $\beta$. \index{regression equation}
\index{least-squares projection}


  There are as many such regressions  as there are ways of selecting
$Y, X$.    We have shown how a model (e.g., a triple $A, C, G$,
together with an initial distribution for $x_0$) restricts a
regression.    Going backward, that is, telling what a given
regression tells about a model, is more difficult. Many regressions tell little about the model, and what little they
have to say can be difficult to decode.    As we indicate in sections \use{sec:estimation1} and \use{sec:estimation2}, the \idx{likelihood
function} completely describes what a given data set says about a model in a way that  is straightforward to decode.

  \subsection{Multiple regressors}
Now let $Y$ be an $n \times 1$ vector of random variables and
think of regression solving the least squares problem for each of
them to attain a representation
$$ Y = \beta X  + \epsilon \EQN leastsq10 $$
where $\beta$ is now $n \times p$ and $\epsilon$ is now an $n
\times 1$ vector of least squares residuals.  The population
regression coefficients are again given by
$$ \beta = E(Y X') [E(X X')]^{-1} . \EQN leastsq12 $$
We will use this formula repeatedly in section \use{sec:Kalman100} to derive
the Kalman filter.\index{Kalman filter}%



